\documentclass[11pt,a4paper]{article}
\usepackage[margin=1in]{geometry}
\usepackage{amsmath,booktabs,hyperref,float,xcolor}
\definecolor{good}{rgb}{0.0,0.5,0.0}

\title{Iteration 017: Combined Loss + Correlation Validation}
\author{Autoresearch Loop}
\date{April 8, 2026}

\begin{document}
\maketitle

\begin{abstract}
Using a combined loss (0.5 MSE + 0.5 correlation) with correlation-based early
stopping achieves a new best $r=0.378$ (+0.5\% over iter011). This fixes the
train/validation metric mismatch identified in iter015 and produces the best
result to date.
\end{abstract}

\section{Hypothesis}
Iter015 showed that correlation loss alone with MSE-based validation selects
suboptimal checkpoints. A combined loss provides stable gradients (from MSE)
while optimizing correlation, and validation on correlation ensures early stopping
selects the correlation-optimal model.

\section{Method}
Loss: $\mathcal{L} = 0.5 \cdot \text{MSE} + 0.5 \cdot (1 - \bar{r})$.
Validation: mean Pearson $r$ across channels (higher is better).
Model: FIR (7 taps) + 15\% channel dropout + CF init. AdamW ($\text{lr}=10^{-3}$),
cosine annealing, 150 epochs.

\section{Results}
\begin{table}[H]
\centering
\begin{tabular}{lrrr}
\toprule
Model & Mean $r$ & Std $r$ & SNR (dB) \\
\midrule
FIR + dropout (iter011) & 0.376 & 0.076 & 0.61 \\
\textbf{Combined loss + corr val} & \textcolor{good}{\textbf{0.378}} & 0.078 & 0.61 \\
\bottomrule
\end{tabular}
\end{table}

\section{Analysis}
The improvement is small (+0.002) but establishes a new best. The combined loss
provides MSE's stable gradients in early training while the correlation component
refines the solution for the actual evaluation metric. The slightly higher std
suggests the model optimizes for per-channel correlation at the expense of some
consistency, which is the correct tradeoff for our metric.

Key insight: matching the training objective to the evaluation metric matters
more than model architecture changes in this regime.

\section{Next}
Iteration 018: Euclidean alignment for cross-subject transfer.

\end{document}
